Para determinar la clase a partir de las observaciones entonces usaremos
\[argmax_y [ P(y_i) \prod_{d}^{x_i} P(x_i | y)]\]
Donde $Y=\{Llueve, Soleado\}$ y $X=(despejado, temperatura\text{ }alta)$ usando Bayes Naïve Frecuentista por lo que tendremos que  aplicar $P(y) = \frac{fr(y)}{N}$, asi que

\begin{align*} % P(c)
    P(Llueve) & = \frac{2}{3} = 0.667\\
    P(Soleado) & = \frac{1}{3} = 0.333 \\\\
\end{align*}

Ahora aplicaremos $P(x|y)=\frac{fr(x,y)}{fr(y)}$ siendo así

\begin{align*} % P(x|c)
    P(despejado|Llueve) = 0\\
    P(temperatura\text{ }alta|Llueve) = \frac{1}{2} = 0.5\\\\
    P(despejado|Soleado) = 1 \\
    P(temperatura\text{ }alta|Soleado) = 1
\end{align*}

Por lo que tenemos lo siguiente para cada clase
\begin{itemize}
    \item Llueve: $P(Llueve, despejado, temperatura\text{ }alta )$
          \begin{align*}
            &= P(despejado|Llueve) \cdot P(temperatura \text{ } alta|Llueve) \cdot P(Llueve)\\
            &= 0 \cdot 0.5 \cdot 0.667\\
            &= 0 
          \end{align*}

    \item Soleado: $ P(Soleado, despejado, temperatura\text{ }alta ) $
        \begin{align*}
           &= P(despejado|Soleado) \cdot P(temperatura \text{ } alta|Soleado) \cdot P(Soleado)\\
            &= 1 \cdot 1 \cdot \frac{1}{3}\\
            &= \frac{1}{3} 
        \end{align*}
\end{itemize}

Ahora que la clase más con el valor más alto es soleado, basta dividir el valor correspondiente a la clase ganadora sobre  los valores obtenidos de cada clase
\[P(Soleado|Observaciones) = \frac{\frac{1}{3}}{\frac{1}{3}+0} = 1\]

Siendo así que la clase más probable es \textbf{Soleado} con $P=1$

 
